\textbf{Phần 1. Trắc nghiệm lựa chọn}
\setcounterpageref{ex}{0}

\begin{ex} %[4D4N17]
$1$ yến bằng bao nhiêu ki-lô-gam?
\choice
{$5\text{ kg}$}
{$100\text{ kg}$}
{$1000\text{ kg}$}
{\True $10\text{ kg}$}
\loigiai{
$1$ yến $= 10\text{ kg}$.\\
\textbf{Đáp án của câu hỏi là D}
}
\end{ex}

\begin{ex} %[4D4N17]
$1$ tạ bằng bao nhiêu ki-lô-gam?
\choice
{$10\text{ kg}$}
{\True $100\text{ kg}$}
{$1000\text{ kg}$}
{$50\text{ kg}$}
\loigiai{
$1$ tạ $= 100\text{ kg}$.\\
\textbf{Đáp án của câu hỏi là B}
}
\end{ex}

\begin{ex} %[4D4N17]
$1$ tấn bằng bao nhiêu ki-lô-gam?
\choice
{$100\text{ kg}$}
{$10\,000\text{ kg}$}
{\True $1000\text{ kg}$}
{$10\text{ kg}$}
\loigiai{
$1$ tấn $= 1000\text{ kg}$.\\
\textbf{Đáp án của câu hỏi là C}
}
\end{ex}

\begin{ex} %[4D4H17]
$1$ tấn bằng bao nhiêu tạ?
\choice
{$5\text{ tạ}$}
{\True $10\text{ tạ}$}
{$100\text{ tạ}$}
{$1000\text{ tạ}$}
\loigiai{
$1$ tấn $= 10\text{ tạ}$.\\
\textbf{Đáp án của câu hỏi là B}
}
\end{ex}

\begin{ex} %[4D4H17]
$3$ tạ $50\text{ kg}$ bằng bao nhiêu ki-lô-gam?
\choice
{$305\text{ kg}$}
{$3500\text{ kg}$}
{$35\text{ kg}$}
{\True $350\text{ kg}$}
\loigiai{
$3 \times 100 + 50 = 350\text{ kg}$.\\
\textbf{Đáp án của câu hỏi là D}
}
\end{ex}

\begin{ex} %[4D4H17]
$2$ tấn $4$ tạ bằng bao nhiêu tạ?
\choice
{\True $24\text{ tạ}$}
{$204\text{ tạ}$}
{$240\text{ tạ}$}
{$6\text{ tạ}$}
\loigiai{
$2 \times 10 + 4 = 24\text{ tạ}$.\\
\textbf{Đáp án của câu hỏi là A}
}
\end{ex}

\begin{ex} %[4D4V17]
Một con voi nặng $3$ tấn $5$ tạ, một con bò nặng $5$ tạ. Con voi nặng gấp mấy lần con bò?
\choice
{$5$ lần}
{$6$ lần}
{$8$ lần}
{\True $7$ lần}
\loigiai{
$3$ tấn $5$ tạ $= 35$ tạ.\\
$35 : 5 = 7$ lần.\\
\textbf{Đáp án của câu hỏi là D}
}
\end{ex}

\begin{ex} %[4D4V17]
Một xe tải chở $40$ bao gạo, mỗi bao nặng $50\text{ kg}$. Xe tải đó chở được bao nhiêu tấn gạo?
\choice
{\True $2\text{ tấn}$}
{$1\text{ tấn}$}
{$3\text{ tấn}$}
{$4\text{ tấn}$}
\loigiai{
Khối lượng: $40 \times 50 = 2000\text{ kg} = 2\text{ tấn}$.\\
\textbf{Đáp án của câu hỏi là A}
}
\end{ex}

\begin{ex} %[4D4V17]
Tính giá trị của biểu thức: $4$ tấn $- 15$ tạ $+ 200\text{ kg}$ theo đơn vị ki-lô-gam.
\choice
{$2500\text{ kg}$}
{$2800\text{ kg}$}
{\True $2700\text{ kg}$}
{$3000\text{ kg}$}
\loigiai{
$4000 - 1500 + 200 = 2700\text{ kg}$.\\
\textbf{Đáp án của câu hỏi là C}
}
\end{ex}

\begin{ex} %[4D4V17]
Một kho có $5$ tấn thóc. Ngày thứ nhất xuất $15$ tạ thóc, ngày thứ hai xuất $20$ tạ thóc. Số thóc còn lại trong kho là
\choice
{$1\text{ tấn } 2\text{ tạ}$}
{\True $1\text{ tấn } 5\text{ tạ}$ ($15\text{ tạ}$)}
{$2\text{ tấn}$}
{$1\text{ tấn } 8\text{ tạ}$}
\loigiai{
$5$ tấn $= 50$ tạ.\\
Còn lại: $50 - (15 + 20) = 15$ tạ $= 1$ tấn $5$ tạ.\\
\textbf{Đáp án của câu hỏi là B}
}
\end{ex}

\begin{ex} %[4D4V17]
Điền dấu thích hợp vào chỗ chấm: $2$ tấn $50\text{ kg} \dots 2500\text{ kg}$.
\choice
{$>$}
{$=$}
{\True $<$}
{Không so sánh được}
\loigiai{
$2$ tấn $50\text{ kg} = 2050\text{ kg} < 2500\text{ kg}$.\\
\textbf{Đáp án của câu hỏi là C}
}
\end{ex}

\begin{ex} %[4D4C17]
Có 3 xe chở hàng: Xe 1 chở $2$ tấn $4$ tạ, xe 2 chở gấp đôi xe 1, xe 3 chở ít hơn trung bình cộng của cả 3 xe là $2$ tạ. Xe 3 chở được số tạ hàng là
\choice
{\True $33$ tạ (hay theo giải hệ trung bình cộng)}
{$32\text{ tạ}$}
{$34\text{ tạ}$}
{$36\text{ tạ}$}
\loigiai{
Xe 1: 24 tạ. Xe 2: 48 tạ.\\
Trung bình cộng 3 xe: $(24 + 48 - 2) : 2 = 35$ tạ.\\
Xe 3 chở: $35 - 2 = 33$ tạ.\\
\textbf{Đáp án của câu hỏi là A}
}
\end{ex}

\textbf{Phần 2. Trắc nghiệm đúng sai}
\setcounterpageref{ex}{0}

\begin{ex} %[4D4H17]
Xét tính đúng sai của các biểu thức đổi đơn vị đo khối lượng:
\choiceTF
{\True $5$ tấn $= 5000\text{ kg}$.}
{$3$ tạ $5\text{ kg} = 350\text{ kg}$.}
{\True $3$ tạ $5\text{ kg} = 305\text{ kg}$.}
{\True $40$ yến $= 4$ tạ.}
\loigiai{
\begin{itemize}
    \item $5$ tấn $= 5000\text{ kg}$. Đúng.
    \item $3$ tạ $5\text{ kg} = 305\text{ kg} \neq 350\text{ kg}$. Ý 2 Sai, Ý 3 Đúng.
    \item $40$ yến $= 400\text{ kg} = 4$ tạ. Đúng.
\end{itemize}
\begin{itemize}
    \item a) \textbf{Đúng}.
    \item b) \textbf{Sai}.
    \item c) \textbf{Đúng}.
    \item d) \textbf{Đúng}.
\end{itemize}
}
\end{ex}

\begin{ex} %[4D4V17]
Hai xe tải chở hàng: Xe A chở $3$ tấn $2$ tạ, xe B chở $28$ tạ. Xét tính đúng sai:
\choiceTF
{\True Xe A chở được $32$ tạ hàng.}
{\True Xe A chở nhiều hơn xe B là $4$ tạ hàng.}
{Cả hai xe chở được tất cả $50$ tạ hàng.}
{\True Cả hai xe chở được tất cả $6$ tấn hàng.}
\loigiai{
\begin{itemize}
    \item Xe A: $32$ tạ. Đúng.
    \item Hiệu: $32 - 28 = 4$ tạ. Đúng.
    \item Tổng: $32 + 28 = 60$ tạ $= 6$ tấn $\neq 50$ tạ. Ý 3 Sai, Ý 4 Đúng.
\end{itemize}
\begin{itemize}
    \item a) \textbf{Đúng}.
    \item b) \textbf{Đúng}.
    \item c) \textbf{Sai}.
    \item d) \textbf{Đúng}.
\end{itemize}
}
\end{ex}

\textbf{Phần 3. Trả lời ngắn}
\setcounterpageref{ex}{0}

\begin{ex} %[4D4H17]
Đổi $4$ tấn $500\text{ kg}$ ra ki-lô-gam.
\shortans{4500}
\loigiai{
\textbf{Đáp số:} $4500$\\
$4500\text{ kg}$.
}
\end{ex}

\begin{ex} %[4D4H17]
Đổi $7$ tạ $6$ yến ra yến.
\shortans{76}
\loigiai{
\textbf{Đáp số:} $76$\\
$7 \times 10 + 6 = 76$ yến.
}
\end{ex}

\begin{ex} %[4D4V17]
Tính tổng: $15$ tạ $+ 25$ tạ theo đơn vị tấn.
\shortans{4}
\loigiai{
\textbf{Đáp số:} $4$\\
$40$ tạ $= 4$ tấn.
}
\end{ex}

\begin{ex} %[4D4V17]
Một con trâu nặng $4$ tạ, một con nghé nặng $120\text{ kg}$. Hỏi con trâu nặng hơn con nghé bao nhiêu ki-lô-gam?
\shortans{280}
\loigiai{
\textbf{Đáp số:} $280$\\
$400 - 120 = 280\text{ kg}$.
}
\end{ex}

\begin{ex} %[4D4V17]
Một chiếc xe chở $6$ bao gạo loại $50\text{ kg}$ và $8$ bao ngô loại $50\text{ kg}$. Xe chở tất cả bao nhiêu tạ lương thực?
\shortans{7}
\loigiai{
\textbf{Đáp số:} $7$\\
$(6 + 8) \times 50 = 700\text{ kg} = 7$ tạ.
}
\end{ex}

\begin{ex} %[4D4C17]
Có hai kho thóc: Kho 1 có $8$ tấn thóc, kho 2 có $6$ tấn thóc. Người ta chuyển từ kho 1 sang kho 2 một số tạ thóc để số thóc ở kho 1 bằng số thóc ở kho 2. Hỏi đã chuyển bao nhiêu tạ thóc?
\shortans{10}
\loigiai{
\textbf{Đáp số:} $10$\\
Hiệu: $8 - 6 = 2$ tấn $= 20$ tạ.\\
Chuyển: $20 : 2 = 10$ tạ.
}
\end{ex}
